\section{Related Work}
\label{sec:related}

\paragraph{Web-agent capability benchmarks.}
Browser-using LLM agents are typically evaluated on what they \emph{can do} under benign conditions: navigation, form-filling, and multi-step task completion. REAL \citep{garg2026real} provides deterministic, locally hosted replicas of eleven popular websites and reports a 41\% top success rate for frontier models on multi-turn tasks, establishing that even capable agents struggle with state-changing web execution. Scammer4U is complementary: we hold capability fixed and ask what the same agents \emph{do under attack}.

\paragraph{Adversarial leakage and dark patterns on the web.}
The closest line of work evaluates web agents against adversarial UI manipulation. \citet{chen2026obviousinvisiblethreat} show that frontier agents indiscriminately process low-salience interface text, with attack success rates up to 74\% for fine-print injections. DECEPTICON \citep{cuvin2026decepticon} reports over 70\% dark-pattern susceptibility across 700 tasks and documents that prompt-level mitigations fail because agents identify the deceptive element in their reasoning trace yet still execute the bad action. TrickyArena \citep{ersoy2026trickyarena} corroborates this on a four-domain testbed of commercial dark patterns, with only a 32\% reduction from prompt-based countermeasures. WebTrap Park \citep{xinyi2026webtrappark} provides a 1{,}226-task automated infrastructure spanning prompt injections and deceptive website designs, including \emph{personal disclosure} as one of its attacker goals. AgentDAM \citep{zharmagambetov2026agentdam} is the closest neighbour by name, measuring \emph{incidental} privacy leakage (12--46\%) on benign sites where the agent over-shares irrelevant PII without any adversary present. Scammer4U is the adversarial counterpart to AgentDAM and the PII-specific counterpart to DECEPTICON: the website itself is the attacker, and an axis-controlled paired-sibling design with pre-registered tests isolates which factors---salience, pressure cues, prompt injection, interaction modality---causally drive critical-PII leakage, evaluated across a four-step mitigation gradient (generic nudge $\rightarrow$ phishing checklist $\rightarrow$ pre-submission reflection) rather than a single countermeasure.

\paragraph{Prompt injection and persuasion attacks.}
AgentDojo \citep{debenedetti2024agentdogo} is the canonical reference for indirect prompt injection, but its 629 security test cases target tool-calling agents on API surfaces (email, banking, travel) rather than agents browsing adversarial websites; targeted attacks succeed in under 25\% of cases there, and the defensive levers (tool filtering, capability gating) differ from web-browsing defenses. TRAP \citep{korgul2025trap} is structurally the most similar to our benchmark, generating attacks from a 5-dimensional matrix that crosses Cialdini persuasion principles with jailbreak templates, and reports a 25\% average ASR. TRAP, however, gates success on a single attacker-controlled redirect click; Scammer4U targets the downstream behaviour after the agent has reached an attacker-controlled form and must decide whether to submit critical PII.

\paragraph{Privacy in multi-agent and social-network settings.}
AgentLeak \citep{yagoubi2026agentleakfullstackbenchmarkprivacy} measures privacy leakage through inter-agent coordination channels, finding that 41.7\% of violations are missed by output-only auditing. AgentSocialBench \citep{wang2026agentsocialbenchevaluatingprivacyrisks} evaluates privacy in agentic social networks and surfaces an \emph{abstraction paradox} whereby prompt-based privacy instructions \emph{increase} partial leakage. \citet{nitu2025machinereadableadsaccessibilitytrust} study how agents engage with commercial ads and dark patterns on a cloned news site, reporting 100\% subscription compliance under goal-contingent pressure. We share the finding that prompt-based defenses are insufficient, but locate the failure on adversarial third-party websites rather than internal coordination, social-graph inference, or commercial nudging.

\paragraph{Phishing corpora and templated environments.}
PhreshPhish \citep{Dalton2025PhreshPhishAR} releases a 17-month real-world phishing-site dataset for passive detection and shows that precision drops sharply at realistic base rates---average precision falls from 0.977 at a 24\% base rate to 0.605 at 0.05\%---while documenting the cloaking, ephemerality, and scrape-noise that make real-world phishing data difficult to use for reproducible study. We use templated sites for this exact reason: scraped corpora lack the interactive backends required for multi-turn form submission and cannot be cleanly axis-controlled for ablation. To our knowledge, no prior benchmark combines (i) attacker-controlled websites that explicitly social-engineer the agent, (ii) axis-controlled paired siblings for single-factor ablation across pressure, salience, and interaction modality, and (iii) a pre-registered analysis plan with falsification criteria for headline claims.

